\documentclass[10pt,a4paper]{article}
% Preamble of the thick-cylinder study report
\usepackage[T1]{fontenc}
\usepackage[utf8]{inputenc}
\usepackage{lmodern}
\usepackage[english]{babel}
\usepackage{microtype}
\usepackage[margin=22mm]{geometry}
\usepackage{amsmath,amssymb,mathtools}
\usepackage{bm}
\usepackage{booktabs}
\usepackage{array}
\usepackage{xcolor}
\usepackage{caption}
\usepackage{subcaption}
\usepackage{tikz}
\usepackage{pgfplots}
\pgfplotsset{compat=1.18}
\usepgfplotslibrary{groupplots}
\usepackage[numbers,sort&compress]{natbib}
\usepackage[colorlinks=true,linkcolor=blue!50!black,citecolor=green!40!black,
            urlcolor=blue!50!black]{hyperref}
\usepackage[nameinlink,capitalise]{cleveref}

\captionsetup{font=small,labelfont=bf}

% data folders (relative to report/)
\newcommand{\tab}{../results/tables}
\newcommand{\pdat}{../mesh/plotdata}

% notation
\newcommand{\PM}{\textsf{PM}}
\newcommand{\PP}{\textsf{PP}}
\newcommand{\PH}{\textsf{PH}}
\newcommand{\eE}{e_{\mathrm{E}}}
\newcommand{\eL}{e_{\mathrm{L}^2}}
\newcommand{\en}{e_{\mathrm{n}}}
\newcommand{\dd}{\,\mathrm{d}}
\newcommand{\bx}{\bm{x}}
\newcommand{\bxi}{\bm{\xi}}

% colours for p = 2..8 (sequential, dark = high p)
\definecolor{p2}{HTML}{FDB863}
\definecolor{p3}{HTML}{F08B3E}
\definecolor{p4}{HTML}{D9541A}
\definecolor{p5}{HTML}{B2182B}
\definecolor{p6}{HTML}{7F1C6A}
\definecolor{p7}{HTML}{4A2A8A}
\definecolor{p8}{HTML}{1B2A5E}

\pgfplotsset{
  conv/.style={
    xmode=log, ymode=log, x dir=reverse,
    width=4.6cm, height=4.6cm,
    xtick={1,0.5,0.25,0.125}, xticklabels={$1$,$\frac12$,$\frac14$,$\frac18$},
    xmin=0.1, xmax=1.25,
    ymin=1e-15, ymax=1e2,
    ytick={1e-14,1e-11,1e-8,1e-5,1e-2,1e1},
    tick label style={font=\scriptsize}, label style={font=\small},
    title style={font=\small, yshift=-1ex},
    unbounded coords=jump, mark size=1.3pt, line width=0.7pt,
    grid=major, grid style={gray!20},
  },
  pline/.style={color=p#1, mark=*, mark options={solid,fill=p#1}},
  pdash/.style={color=p#1, mark=o, dashed, mark options={solid}},
  phyb/.style={color=p#1, mark=square*, dashdotted, mark options={solid, fill=p#1}},
  meshplot/.style={
    width=3.35cm, height=3.35cm, axis equal image, hide axis,
    xmin=-0.05, xmax=3.05, ymin=-0.05, ymax=3.05,
    title style={font=\scriptsize, yshift=-1.2ex},
  },
}

% the same styles for TikZ paths (legends drawn outside the axes)
\tikzset{
  pline/.style={color=p#1, line width=0.7pt},
  pdash/.style={color=p#1, dashed, line width=0.7pt},
  phyb/.style={color=p#1, dashdotted, line width=0.7pt},
}

% h-convergence panel: #1 problem, #2 form, #3 alpha (3 digits), #4 key (eE, eL, eN)
\newcommand{\hpanel}[4]{%
  \foreach \pp in {2,...,8}{%
    \edef\tmp{\noexpand\addplot[pline=\pp] table[x=h, y=#4\pp] {\tab/h_#1_#2_a#3.dat};}\tmp
  }%
}
% slope triangle: #1 x-start (h), #2 y at x-start, #3 slope
\newcommand{\slope}[3]{%
  \draw[gray, thin] (axis cs:#1,#2) -- (axis cs:#1/2,#2/2^#3) -- (axis cs:#1,#2/2^#3) -- cycle;
  \node[gray, font=\tiny, anchor=west] at (axis cs:#1,{#2/2^(#3/2)}) {#3};
}

\usetikzlibrary{calc}
\renewcommand{\tab}{../results}

\title{Reproduction of the nodal rotation patch test and the convergence test of
       Rajendran (2010) for $\theta\ge90^\circ$ with the USEM code base}
\author{USEM-MATLAB study \texttt{studies/rajendran\_panel}}
\date{\today}

\begin{document}
\maketitle

\begin{abstract}
\noindent
The square-panel tests of \citet[Secs.~5.5--5.7]{Rajendran2010technique} (mid-side nodes
rotated by $\theta=90^\circ$--$270^\circ$, overlapping and folded elements) are recomputed
with the spectral elements of the code base on gmsh meshes. With the 9-node element defined
as in the paper (parametric test functions, metric trial functions in Cartesian coordinates,
9th node at the centroid of the corners, $5\times5$ Gauss points), all values of the 9-node
Tables~15--19 are reproduced to four digits, for the unsymmetric (PM-QUAD9$^\ast$) as well as
the conventional (PP-QUAD9) and the metric (MM-QUAD9) elements, except for entries that are
inconsistent within the paper itself (\cref{sec:paper}). The energy-norm convergence rates of
PM-QUAD9$^\ast$ (Fig.~15 of the paper) are reproduced as well ($2.00$ and $1.98$ against
$2.11$ and $1.97$). For $p\ge3$ the unsymmetric spectral element reproduces the cubic field of
the shear test on every mesh, including fully folded ones, up to round-off. The code base as
is, however, places interior nodes with the 8-node map; on these meshes this makes interior
and edge nodes coincide for even $p$, and the solution fails (\cref{sec:code}).
\end{abstract}

\section{Set-up}\label{sec:setup}
Panel $[0,2]\times[-1,1]$, plane stress, $\nu=0.25$, three point supports ($u_x=0$ at
$(0,\pm1)$, $u_x=u_y=0$ at $(0,0)$), equilibrated loads on $x=0$ and $x=2$ (tension
$\sigma_x=1000$; moment $\sigma_x=-1500y$; shear $\sigma_x=-1500(2-x)y$,
$\tau_{xy}=750(1-y^2)$), displacement at $A=(2,0)$. The meshes (\cref{fig:meshes}) are
generated with gmsh (\texttt{mesh/gen\_panel\_meshes.py}): $2\times2$ super-elements of
$2\times2$ elements; the mid-side nodes of the internal edges of every super-element are
rotated about its centre by $\theta$ (Meshes a--f: $\theta=0^\circ,90^\circ,135^\circ,175^\circ,
225^\circ,270^\circ$; sequences A, B, C: $\theta=0^\circ,90^\circ,135^\circ$ with $1\times1$,
$2\times2$, $4\times4$ super-elements). For $\theta\ge90^\circ$ every element of every mesh has
$\det\bm J\le0$ somewhere.

\paragraph{Material constant.} The paper states $E=1500$, but its exact solutions ($1.0$,
$1.5$, $3.3125$) are those of $E=2000$ (elasticity solutions with the supports above: $1.333$,
$2.0$, $4.4167$ for $E=1500$; $1.0$, $1.5$, $3.3125$ for $E=2000$). $E=2000$ is used here; all
displacements scale with $1/E$.

\paragraph{Loads.} The consistent nodal forces (\texttt{matlab/rp\_loads.m}) agree with the
lists of the paper (Fig.~8 and Sec.~5.6) except for the moment case of Fig.~8, where the
paper lists $-250$, $-500$ at nodes 8, 9 (and $250$, $500$ at nodes 14, 13); the consistent,
antisymmetric values are $-500$, $-250$ (and $500$, $250$). The listed values would give a
moment of $1125$ instead of $1000$; since the paper's PM results equal the exact $1.5$, the
consistent values were evidently used.

\paragraph{Elements.} With $p=2$ (9 nodes):
PP9 = conventional element; MM9 = metric functions on both sides, $\int\bm B_R^{\mathsf T}\bm D
\bm B_R\det\bm J$; PM9 = PM-QUAD9$^\ast$, $\int(\bm B_L^\ast)^{\mathsf T}\bm D\bm B_R$. As in the
paper: metric functions $Q_2$ in Cartesian coordinates relative to the element centroid
(\texttt{ShapeFunctions2D\_Metric}, method 0), 9th node at the centroid of the corners,
$5\times5$ Gauss points. In addition: ``\PM, skew'' = the code-base metric functions (skew
coordinates, method 1) with the same nodes; ``\PM, code'' = \texttt{Ke\_2D} with the nodes of
\texttt{pSpectralElement\_NodesTopology} ($4\times4$ points). For $p=3,\dots,8$: \PP\ and \PM\
with interior nodes from the bilinear corner map (the generalisation of the corner centroid)
and $2p\times2p$ points, and \PM$_\mathrm{code}$ with the interior nodes of the code base.

\begin{figure}[t]
  \centering
  \begin{tikzpicture}
\begin{groupplot}[group style={group size=3 by 2, horizontal sep=0.05\textwidth, vertical sep=7mm}, scale only axis, width=0.3\textwidth, height=0.3\textwidth, hide axis, xmin=-0.03, xmax=2.03, ymin=-1.03, ymax=1.03, title style={font=\small, yshift=-0.5ex}]
\nextgroupplot[title={Mesh a: $\theta=0^\circ$}]
\addplot[black, line width=0.4pt] table {../mesh/plotdata/mesh_N1_t000.dat};
\nextgroupplot[title={Mesh b: $\theta=90^\circ$}]
\addplot[black, line width=0.4pt] table {../mesh/plotdata/mesh_N1_t090.dat};
\nextgroupplot[title={Mesh c: $\theta=135^\circ$}]
\addplot[black, line width=0.4pt] table {../mesh/plotdata/mesh_N1_t135.dat};
\nextgroupplot[title={Mesh d: $\theta=175^\circ$}]
\addplot[black, line width=0.4pt] table {../mesh/plotdata/mesh_N1_t175.dat};
\nextgroupplot[title={Mesh e: $\theta=225^\circ$}]
\addplot[black, line width=0.4pt] table {../mesh/plotdata/mesh_N1_t225.dat};
\nextgroupplot[title={Mesh f: $\theta=270^\circ$}]
\addplot[black, line width=0.4pt] table {../mesh/plotdata/mesh_N1_t270.dat};
\end{groupplot}
\end{tikzpicture}

\medskip

\begin{tikzpicture}
\begin{groupplot}[group style={group size=4 by 1, horizontal sep=0.04\textwidth, vertical sep=7mm}, scale only axis, width=0.22\textwidth, height=0.22\textwidth, hide axis, xmin=-0.03, xmax=2.03, ymin=-1.03, ymax=1.03, title style={font=\small, yshift=-0.5ex}]
\nextgroupplot[title={B, $2\times2$}]
\addplot[black, line width=0.4pt] table {../mesh/plotdata/mesh_N2_t090.dat};
\nextgroupplot[title={B, $4\times4$}]
\addplot[black, line width=0.4pt] table {../mesh/plotdata/mesh_N4_t090.dat};
\nextgroupplot[title={C, $2\times2$}]
\addplot[black, line width=0.4pt] table {../mesh/plotdata/mesh_N2_t135.dat};
\nextgroupplot[title={C, $4\times4$}]
\addplot[black, line width=0.4pt] table {../mesh/plotdata/mesh_N4_t135.dat};
\end{groupplot}
\end{tikzpicture}

  \caption{gmsh meshes: Meshes a--f of the patch test and the $2\times2$, $4\times4$ meshes of
    sequences B and C (cf.\ Figs.~9, 10 of the paper).}
  \label{fig:meshes}
\end{figure}

\section{Nodal rotation patch test ($p=2$)}\label{sec:patch}
\begin{table}[h]
  \centering\small
  \caption{Moment loading, $u_y(A)$ (paper: Table~15; exact $1.5$).}
  \label{tab:moment}
  \resizebox{\linewidth}{!}{\begin{tabular}{l rr rr rr rr}
\toprule
 & \multicolumn{2}{c}{PP9} & \multicolumn{2}{c}{MM9} & \multicolumn{2}{c}{PM9 (PM-QUAD9$^\ast$)} & \multicolumn{2}{c}{\PM, $p=2$}\\
\cmidrule(lr){2-3}\cmidrule(lr){4-5}\cmidrule(lr){6-7}\cmidrule(lr){8-9}
mesh & present & paper & present & paper & present & paper & skew & code\\
\midrule
a (0$^\circ$) & 1.5000 & 1.5000 & 1.5000 & 1.5000 & 1.5000 & 1.5000 & 1.5000 & 1.5000\\
b (90$^\circ$) & 1.2606 & 1.2606 & 1.7935 & 1.7935 & 1.5000 & 1.5000 & 1.5000 & err.\\
c (135$^\circ$) & 1.0372 & 1.0372 & 1.6398 & 1.6398 & 1.5000 & 1.5000 & 1.5000 & 1.5000\\
d (175$^\circ$) & 1.2472 & 1.2472 & 1.5285 & 1.5285 & 1.5000 & 1.5000 & 1.5000 & 1.4996\\
e (225$^\circ$) & 1.0372 & 1.0372 & 1.6398 & 1.6398 & 1.5000 & 1.5000 & 1.5000 & 1.5000\\
f (270$^\circ$) & 1.2606 & 1.2606 & 1.7935 & 1.7935 & 1.5000 & 1.5000 & 1.5000 & err.\\
\bottomrule
\end{tabular}
}

  \vspace{3mm}
  \caption{Shear loading, $u_y(A)$ (paper: Table~16; exact $3.3125$). $^\ast$: differs from the
    present result, see \cref{sec:paper}.}
  \label{tab:shear}
  \resizebox{\linewidth}{!}{\begin{tabular}{l rr rr rr rr}
\toprule
 & \multicolumn{2}{c}{PP9} & \multicolumn{2}{c}{MM9} & \multicolumn{2}{c}{PM9 (PM-QUAD9$^\ast$)} & \multicolumn{2}{c}{\PM, $p=2$}\\
\cmidrule(lr){2-3}\cmidrule(lr){4-5}\cmidrule(lr){6-7}\cmidrule(lr){8-9}
mesh & present & paper & present & paper & present & paper & skew & code\\
\midrule
a (0$^\circ$) & 3.2979 & 3.2979 & 3.2979 & 3.2979 & 3.2979 & 3.2979 & 3.2979 & 3.2979\\
b (90$^\circ$) & 3.0033 & 1.9088$^\ast$ & 3.9559 & 3.9559 & 3.2710 & 3.2710 & 3.3120 & err.\\
c (135$^\circ$) & -0.4993 & -0.4993 & 3.8294 & 3.8294 & 3.3044 & 3.3044 & 3.4134 & 4.1224\\
d (175$^\circ$) & 2.5498 & 2.5498 & 3.3657 & 3.3657 & 3.3677 & 3.3677 & 3.1501 & 3.0340\\
e (225$^\circ$) & -0.4993 & -0.4993 & 3.8294 & 3.8294 & 3.3044 & 3.3044 & 3.4134 & 4.1224\\
f (270$^\circ$) & 3.0033 & 1.9088$^\ast$ & 3.9559 & 3.9559 & 3.2710 & 3.2710 & 3.3120 & err.\\
\bottomrule
\end{tabular}
}

  \vspace{3mm}
  \caption{Tension, $u_x(A)$ (paper: Table~14; exact $1.0$).}
  \label{tab:tension}
  \resizebox{\linewidth}{!}{\begin{tabular}{l rr rr rr rr rr}
\toprule
 & \multicolumn{2}{c}{PP9} & \multicolumn{2}{c}{MM9} & \multicolumn{2}{c}{PM9 (PM-QUAD9$^\ast$)} & \multicolumn{2}{c}{\PM, $p=2$} & \multicolumn{2}{c}{\PH, $p=2$}\\
\cmidrule(lr){2-3}\cmidrule(lr){4-5}\cmidrule(lr){6-7}\cmidrule(lr){8-9}\cmidrule(lr){10-11}
mesh & present & paper & present & paper & present & paper & skew & code & centr. & code\\
\midrule
a (0$^\circ$) & $1.0000$ & $1.0$ & $1.0000$ & $1.0$ & $1.0000$ & $1.0$ & $1.0000$ & $1.0000$ & $1.0000$ & $1.0000$\\
b (90$^\circ$) & $1.0000$ & $1.0$ & $1.1300$ & $1.0^\ast$ & $1.0000$ & $1.0$ & $1.0000$ & err. & $1.0000$ & err.\\
c (135$^\circ$) & $1.0000$ & $1.0$ & $1.4633$ & $1.0^\ast$ & $1.0000$ & $1.0$ & $1.0000$ & $1.0000$ & $1.0000$ & $1.0000$\\
d (175$^\circ$) & $1.0000$ & $1.0$ & $1.0285$ & $1.0^\ast$ & $1.0000$ & $1.0$ & $1.0000$ & $1.0000$ & $1.0000$ & $1.0000$\\
e (225$^\circ$) & $1.0000$ & $1.0$ & $1.4633$ & $1.0^\ast$ & $1.0000$ & $1.0$ & $1.0000$ & $1.0000$ & $1.0000$ & $1.0000$\\
f (270$^\circ$) & $1.0000$ & $1.0$ & $1.1300$ & $1.0^\ast$ & $1.0000$ & $1.0$ & $1.0000$ & err. & $1.0000$ & err.\\
\bottomrule
\end{tabular}
}
\end{table}
PM9 passes the constant and the linear strain patch tests on all meshes, as in the paper
(\cref{tab:tension,tab:moment}); for the quadratic-strain (cubic displacement) shear test, which
lies outside $Q_2$, its results are within $1.7\,\%$ of the exact value on all meshes
(\cref{tab:shear}). All PP9 and MM9 values for moment and shear coincide with the paper to four
digits. The metric functions of the code base (skew coordinates) give the same results for the
constant and linear strain tests and comparable ones for the shear test ($3.15$--$3.41$).

\section{Convergence test (shear loading, $p=2$)}\label{sec:conv}
\begin{table}[h]
  \centering\small
  \caption{$u_y(A)$ for the sequences A, B, C (paper: Tables~17--19; exact $3.3125$).}
  \label{tab:conv}
  \resizebox{\linewidth}{!}{\begin{tabular}{ll rr rr rr rr rr}
\toprule
 & & \multicolumn{2}{c}{PP9} & \multicolumn{2}{c}{MM9} & \multicolumn{2}{c}{PM9} & \multicolumn{2}{c}{\PM, $p=2$} & \multicolumn{2}{c}{\PH, $p=2$}\\
\cmidrule(lr){3-4}\cmidrule(lr){5-6}\cmidrule(lr){7-8}\cmidrule(lr){9-10}\cmidrule(lr){11-12}
seq. & mesh & present & paper & present & paper & present & paper & skew & code & centr. & code\\
\midrule
A & $1\times1$ & 3.2979 & 3.2979 & 3.2979 & 3.2979 & 3.2979 & 3.2979 & 3.2979 & 3.2979 & 3.2979 & 3.2979\\
 & $2\times2$ & 3.3106 & 3.3106 & 3.3106 & 3.3106 & 3.3106 & 3.3106 & 3.3106 & 3.3106 & 3.3106 & 3.3106\\
 & $4\times4$ & 3.3123 & 3.3123 & 3.3123 & 3.3123 & 3.3123 & 3.3123 & 3.3123 & 3.3123 & 3.3123 & 3.3123\\
\midrule
B & $1\times1$ & 3.0033 & 3.0033 & 3.9559 & 3.9559 & 3.2710 & 3.3044$^\ast$ & 3.3120 & err. & 3.3723 & err.\\
 & $2\times2$ & 3.1768 & 3.1768 & 3.9519 & 3.9519 & 3.3000 & 3.3057$^\ast$ & 3.3185 & err. & 3.3229 & err.\\
 & $4\times4$ & 3.2662 & 3.2661 & 3.9598 & 3.9598 & 3.3110 & 3.3116$^\ast$ & 3.3121 & err. & 3.3149 & err.\\
\midrule
C & $1\times1$ & -0.4993 & -0.4993 & 3.8294 & 3.8294 & 3.3044 & 3.2710$^\ast$ & 3.4134 & 4.1224 & 3.4016 & 3.0874\\
 & $2\times2$ & 2.9218 & 2.9218 & 3.9208 & 3.9208 & 3.3057 & 3.3001$^\ast$ & 3.2600 & 3.2935 & 3.3218 & 3.2938\\
 & $4\times4$ & 3.1531 & 3.1531 & 3.9660 & 3.9660 & 3.3116 & 3.3110$^\ast$ & 3.3067 & 3.3105 & 3.3121 & 3.3075\\
\bottomrule
\end{tabular}
}

  \vspace{3mm}
  \caption{Least-squares rates of the energy-norm strain error, Eq.~(24) of the paper (paper:
    Fig.~15). \PP\ on folded meshes: the parametric strain is singular where $\det\bm J=0$, so the
    energy error of \PP\ is not well defined for B and C (here: signed $\det\bm J$).
    ``--'': error at round-off level (field reproduced).}
  \label{tab:rates}
  \begin{tabular}{l rr rr rrrr}
\toprule
 & \multicolumn{2}{c}{PP9} & \multicolumn{2}{c}{PM9} & \PM, $p=2$ & \PM, $p=2$ & \PP, $p=3$ & \PM, $p=3$\\
\cmidrule(lr){2-3}\cmidrule(lr){4-5}
seq. & present & paper & present & paper & skew & code & & \\
\midrule
A & 1.99 & 1.93 & 1.99 & 1.93 & 1.99 & 1.99 & -- & --\\
B & 0.70 & 1.05 & 2.00 & 2.11 & 2.02 & 0.01 & 0.64 & --\\
C & 2.50 & 0.85 & 1.98 & 1.97 & 2.13 & 2.41 & 0.94 & --\\
\bottomrule
\end{tabular}

\end{table}
The displacements of sequence A and of PP9 and MM9 for B and C agree with the paper. PM9 on
sequence B yields the values the paper lists for sequence C and vice versa (\cref{sec:paper}).
PM9 converges with the optimal rate $2$ on the folded meshes of B and C, as reported in the paper.
For $p=3$ the unsymmetric element reproduces the field on all meshes (error $\le2\cdot10^{-12}$),
whereas \PP\ ($p=3$) still has energy-norm errors of $9$--$80\,\%$.

\section{Higher degrees}\label{sec:highp}
\begin{table}[h]
  \centering\small
  \caption{Shear loading on Meshes a--f: relative nodal error $e_n$ for $p=3,\dots,8$ (cubic
    field, contained in the \PM\ trial space for $p\ge3$). \PM: interior nodes from the bilinear
    corner map; \PM$_\mathrm{code}$: interior nodes of \texttt{pSpectralElement\_NodesTopology}.}
  \label{tab:highp}
  \begin{tabular}{lrrrrrr}
\toprule
 & a (0$^\circ$) & b (90$^\circ$) & c (135$^\circ$) & d (175$^\circ$) & e (225$^\circ$) & f (270$^\circ$)\\
\midrule
\PP, $p=3$ & 1e-14 & 2e-1 & 1e-1 & 2e-1 & 1e-1 & 2e-1\\
\PM, $p=3$ & 2e-14 & 2e-14 & 7e-14 & 1e-13 & 2e-14 & 3e-14\\
\PM$_\mathrm{code}$, $p=3$ & 7e-14 & 7e-13 & 1e-12 & 7e-8 & 6e-12 & 1e-13\\
\PH, $p=3$ & 3e-14 & 3e-14 & 3e-13 & 2e-13 & 2e-13 & 6e-13\\
\PH$_\mathrm{code}$, $p=3$ & 2e-14 & 5e-14 & 9e-14 & 6e-8 & 8e-14 & 1e-14\\
\addlinespace[2pt]
\PP, $p=4$ & 1e-13 & 8e-2 & 5e-1 & 2e-1 & 5e-1 & 8e-2\\
\PM, $p=4$ & 1e-13 & 2e-13 & 3e-12 & 3e-11 & 3e-12 & 2e-13\\
\PM$_\mathrm{code}$, $p=4$ & 1e-13 & err. & 3e-11 & err. & 5e-10 & err.\\
\PH, $p=4$ & 1e-13 & 2e-14 & 6e-12 & 5e-13 & 1e-11 & 1e-14\\
\PH$_\mathrm{code}$, $p=4$ & 1e-13 & err. & 1e-12 & err. & 1e-12 & err.\\
\addlinespace[2pt]
\PP, $p=5$ & 8e-13 & 7e-2 & 9e-2 & 2e-1 & 9e-2 & 7e-2\\
\PM, $p=5$ & 9e-13 & 1e-10 & 6e-13 & 8e-11 & 6e-13 & 3e-11\\
\PM$_\mathrm{code}$, $p=5$ & 8e-13 & 4e-10 & 8e-10 & err. & 6e-10 & 5e-9\\
\PH, $p=5$ & 9e-13 & 8e-13 & 2e-12 & 2e-11 & 5e-12 & 2e-13\\
\PH$_\mathrm{code}$, $p=5$ & 6e-13 & 7e-13 & 2e-12 & err. & 2e-12 & 2e-13\\
\addlinespace[2pt]
\PP, $p=6$ & 7e-13 & 9e-2 & 2e-1 & 1e-1 & 2e-1 & 9e-2\\
\PM, $p=6$ & 6e-13 & 1e-11 & 6e-11 & 3e-11 & 2e-11 & 1e-11\\
\PM$_\mathrm{code}$, $p=6$ & 6e-13 & err. & 1e-8 & err. & 7e-9 & err.\\
\PH, $p=6$ & 5e-13 & 2e-13 & 2e-11 & 1e-10 & 1e-11 & 2e-13\\
\PH$_\mathrm{code}$, $p=6$ & 4e-13 & err. & 2e-12 & err. & 2e-12 & err.\\
\addlinespace[2pt]
\PP, $p=7$ & 3e-13 & 1e-1 & 9e-2 & 1e-1 & 9e-2 & 1e-1\\
\PM, $p=7$ & 4e-14 & 2e-11 & 4e-10 & 8e-10 & 1e-9 & 4e-11\\
\PM$_\mathrm{code}$, $p=7$ & 3e-13 & 2e-6 & 8e-8 & err. & 4e-7 & 2e-6\\
\PH, $p=7$ & 5e-13 & 5e-12 & 1e-10 & 3e-10 & 2e-11 & 9e-12\\
\PH$_\mathrm{code}$, $p=7$ & 1e-12 & 2e-12 & 4e-12 & err. & 3e-11 & 2e-12\\
\addlinespace[2pt]
\PP, $p=8$ & 3e-12 & 1e-1 & 1e0 & 1e-1 & 1e0 & 1e-1\\
\PM, $p=8$ & 3e-12 & 3e-9 & 8e-6 & 1e-7 & 1e-5 & 5e-9\\
\PM$_\mathrm{code}$, $p=8$ & 3e-12 & err. & err. & err. & err. & err.\\
\PH, $p=8$ & 3e-12 & 1e-10 & 7e-8 & 6e-9 & 1e-8 & 4e-11\\
\PH$_\mathrm{code}$, $p=8$ & 2e-12 & err. & 1e-11 & err. & 3e-11 & err.\\
\bottomrule
\end{tabular}

\end{table}
With interior nodes from the bilinear corner map, the unsymmetric element is Jacobian-unlimited
immune for $p=3,\dots,8$ (\cref{tab:highp}): the cubic field is reproduced on all meshes with
overlapping and folded elements. The round-off level grows with $p$ and distortion
(from $10^{-13}$ to $10^{-5}$ for $p=8$, Mesh c), in line with the thick-cylinder study. The
conventional element has errors of $7$--$100\,\%$ on the distorted meshes for all $p$.

\section{A defect of the code base on these meshes}\label{sec:code}
\texttt{pSpectralElement\_NodesTopology} places the interior nodes with the 8-node (serendipity)
map. On Meshes b and f this maps the interior node $\bxi=(0,0)$ of the $p=2$ element exactly onto
a mid-side node of the same element (e.g.\ $(1,-0.5)$); the metric Vandermonde matrix is
singular ($\mathrm{cond}\approx10^{18}$--$10^{20}$) and the computed displacements are zero
(``code'' columns of \cref{tab:moment,tab:shear,tab:tension,tab:conv}). The same happens for every
even $p$ (interior nodes on $\xi=0$ or $\eta=0$), and on Meshes c--e the near-coincidence of
nodes destroys the accuracy for $p\ge4$ (\cref{tab:highp}). Placing the interior nodes with the
bilinear map of the corners, which for $p=2$ is the choice of the paper (``for simplicity''),
removes the problem on all meshes tested. \texttt{ShapeFunctions2D\_Metric} uses
\texttt{decomposition}, which does not stop on a singular matrix, so the failure is silent.

\section{Inconsistencies in the paper}\label{sec:paper}
\begin{enumerate}
\item $E=1500$ is stated, the exact values correspond to $E=2000$ (\cref{sec:setup}).
\item Fig.~8, moment loading: nodal forces at nodes 8/9 and 14/13 interchanged (\cref{sec:setup}).
\item Table~16, PP-QUAD9, Meshes b and f: $1.9088$; the present value $3.0033$ equals the paper's
  own PP-QUAD9 value for the same mesh in Table~18 ($1\times1$ super-element of sequence B).
\item Tables~18 and 19, PM-QUAD9$^\ast$ columns: interchanged. The present values for sequence B
  ($3.2710$, $3.3000$, $3.3110$) are the paper's Table~19 entries ($3.2710$, $3.3001$, $3.3110$),
  those for C the Table~18 entries; the $1\times1$ values of Table~16 (b: $3.2710$, c: $3.3044$)
  confirm the present assignment.
\item Table~14, MM-QUAD9: $1.0$ for all meshes. With a stiffness matrix that reproduces all MM9
  moment and shear values to four digits, MM9 gives $1.1300$ (b, f), $1.4633$ (c, e), $1.0285$ (d),
  also with load vectors computed from the metric functions (Eq.~(22) of the paper). MM9 does not
  pass the constant strain patch test on these meshes; the paper's entries could not be reproduced.
\end{enumerate}
Apart from these items, every 9-node value of the paper for $\theta\ge90^\circ$ is reproduced to
four digits. The 8-node results (Tables~8--13) are not reproduced: the code base provides only
tensor-product (Lagrange) elements.

\section{Reproduction}
\texttt{mesh/gen\_panel\_meshes.py} (gmsh), \texttt{matlab/rp\_run.m} (all computations, GNU
Octave 8.4 or MATLAB, a few minutes), \texttt{postprocess.py} (tables). Paper values:
\texttt{results/paper\_values.csv}, \texttt{results/paper\_rates.csv}.

\bibliographystyle{plainnat}
\bibliography{../../../literature/references}
\end{document}
